\section{A finite sieve theorem and effective asymptotic Brun--Titchmarsh}
\label{sec:finite-sieve}

This section proves a rational inequality that already has the correct
uniform dependence on the modulus.  It also proves a version of
Brun--Titchmarsh with constant $2+1/k$ and an explicit primitive recursive
threshold.  The sharp constant $2$ for every $x>q$ is a stronger assertion;
the distinction will be maintained throughout.

For natural numbers $x,a$ and a positive integer $q$, define
\[
 \pi(x;q,a)=\#\{p\leq x:p\text{ is prime and }p\equiv a\pmod q\}.
\]
For $z\geq1$ put
\begin{align*}
 D_q(z)&=\{d:1\leq d\leq z,\ d\text{ is squarefree},\ \gcd(d,q)=1\},\\
 G_q(z)&=\sum_{d\in D_q(z)}\frac1{\varphi(d)},
 & H_z&=\sum_{n=1}^{z}\frac1n.
\end{align*}
Here $\varphi(1)=1$, and $\mu$ denotes the M\"obius function, defined by
bounded prime factorization.  In particular, every displayed sum in these
definitions is a finite rational number, and $G_q(z),H_z>0$ because their
$d=1$ and $n=1$ summands are present.

\begin{theorem}[Finite rational sieve inequality]
\label{thm:finite-sieve}
If $q,z\geq1$ and $\gcd(a,q)=1$, then
\begin{equation}
 \boxed{\displaystyle
 \pi(x;q,a)\leq \frac{x}{qG_q(z)}+z^2+z
 \leq \frac{x}{\varphi(q)H_z}+z^2+z.}
 \label{eq:finite-sieve}
\end{equation}
\end{theorem}

The statement becomes a quantifier-free integer inequality after
clearing its positive denominators.  We give all the identities needed
for its proof; in particular, we will not use the existence of a
minimizer for a quadratic form.

\subsection{Elementary divisor identities}

We use
\begin{equation}
 \sum_{r\mid n}\varphi(r)=n,\qquad
 \sum_{r\mid n}\mu(r)=
 \begin{cases}1&n=1,\\0&n>1.\end{cases}
 \label{eq:divisor-identities}
\end{equation}
The first identity follows by partitioning $1,\ldots,n$ according to
their greatest common divisor with $n$: the part with greatest common
divisor $n/r$ has $\varphi(r)$ elements.  For the second, if $n>1$, let
$p$ be its least prime divisor.  Pair the squarefree divisors not
containing $p$ with the divisors obtained by adjoining $p$; their
M\"obius signs cancel.  Divisors that are not squarefree contribute
zero.  Both arguments are finite.

The Chinese remainder theorem gives
$\varphi(uv)=\varphi(u)\varphi(v)$ whenever $\gcd(u,v)=1$.
For squarefree $d$, expanding a finite product therefore gives
\begin{equation}
 \sum_{e\mid d}\frac1{\varphi(e)}
 =\prod_{p\mid d}\left(1+\frac1{p-1}\right)
 =\frac d{\varphi(d)}.
 \label{eq:phi-inverse-divisors}
\end{equation}
All prime and divisor lists here can be generated by trial division
bounded by the relevant integer.

\subsection{Writing down the sieve weights}

Write $D=D_q(z)$ and $G=G_q(z)$.  The set $D$ is closed under taking
positive divisors.  For a rational vector $(\lambda_d)_{d\in D}$ set
\[
 y_r=\sum_{\substack{d\in D\\r\mid d}}\frac{\lambda_d}{d}.
\]
The identities $[d,e]=de/\gcd(d,e)$ and
\eqref{eq:divisor-identities} give the exact diagonalization
\begin{equation}
 Q(\lambda):=\sum_{d,e\in D}\frac{\lambda_d\lambda_e}{[d,e]}
 =\sum_{r\in D}\varphi(r)y_r^2.
 \label{eq:sieve-diagonal}
\end{equation}
Indeed, substitute
$\gcd(d,e)=\sum_{r\mid d,\,r\mid e}\varphi(r)$ into the double sum and
interchange three finite sums.

Finite M\"obius inversion recovers the weights from $y$:
\begin{equation}
 \lambda_d=d\sum_{\substack{r\in D\\d\mid r}}\mu(r/d)y_r.
 \label{eq:sieve-inversion}
\end{equation}
To check the formula, substitute it into the defining sum for $y_s$.
The coefficient of $y_r$ becomes
$\sum_{s\mid d\mid r}\mu(r/d)$, which equals one when $r=s$ and zero
otherwise.  All intermediate divisors belong to $D$.

Now make the explicit rational choice
\begin{equation}
 y_r=\frac{\mu(r)}{\varphi(r)G},
 \qquad
 \lambda_d=d\sum_{\substack{r\in D\\d\mid r}}
               \frac{\mu(r/d)\mu(r)}{\varphi(r)G}.
 \label{eq:explicit-sieve-weights}
\end{equation}
Since $\mu(r)^2=1$ on $D$, equations
\eqref{eq:sieve-diagonal} and \eqref{eq:sieve-inversion} show that
\begin{equation}
 \lambda_1=1,\qquad Q(\lambda)=\frac1G.
 \label{eq:sieve-normalization}
\end{equation}
These weights are in fact optimal for the quadratic form under the
constraint $\lambda_1=1$.  The following identity proves that assertion
without an appeal to analysis:
\[
 Q(\lambda)-\frac1G=
 \sum_{r\in D}\varphi(r)
 \left(y_r-\frac{\mu(r)}{\varphi(r)G}\right)^2\geq0.
\]
To expand the right side, use
$\sum_{r\in D}\mu(r)y_r=\lambda_1=1$.

For later use we prove a uniform bound on the weights.  Define
\[
 G_{dq}(z/d)=
 \sum_{\substack{s\leq z/d\\s\text{ squarefree}\\\gcd(s,dq)=1}}
 \frac1{\varphi(s)}.
\]
Since the indices in \eqref{eq:explicit-sieve-weights} are squarefree,
writing $r=ds$ yields
\begin{equation}
 \lambda_d=\frac{\mu(d)d}{\varphi(d)G}G_{dq}(z/d).
 \label{eq:weight-formula}
\end{equation}
If $e\mid d$ and $s$ occurs in $G_{dq}(z/d)$, then $es\in D$.
The map $(e,s)\mapsto es$ is injective: $e$ can be recovered as
$\gcd(es,d)$ and then $s$ can be recovered by division.  Therefore
\[
 G\geq\sum_{e\mid d}\frac1{\varphi(e)}G_{dq}(z/d)
 =\frac d{\varphi(d)}G_{dq}(z/d).
\]
Together with \eqref{eq:weight-formula}, this gives
\begin{equation}
 |\lambda_d|\leq1,
 \qquad \sum_{d\in D}|\lambda_d|\leq |D|\leq z.
 \label{eq:weight-bound}
\end{equation}

\subsection{Applying the quadratic majorant}

Let $P$ be the product of the primes $p\leq z$ that do not divide $q$.
For every positive integer $n$,
\begin{equation}
 \boldsymbol{1}_{\gcd(n,P)=1}
 \leq\left(\sum_{\substack{d\in D\\d\mid n}}\lambda_d\right)^2.
 \label{eq:sieve-majorant}
\end{equation}
If $\gcd(n,P)=1$, only $d=1$ occurs, so both sides equal one.
Otherwise the left side is zero and the right side is a rational
square.

For $l=[d,e]$, we have $\gcd(l,q)=1$.  The simultaneous conditions
$n\equiv a\pmod q$ and $l\mid n$ select one residue class modulo $ql$.
Counting that class in $1,\ldots,x$ gives
\begin{equation}
 \#\{1\leq n\leq x:n\equiv a\pmod q,\ l\mid n\}
 =\frac{x}{ql}+r_l,
 \qquad |r_l|\leq1.
 \label{eq:crt-count}
\end{equation}
For example, choose its representative in $\{1,\ldots,ql\}$ and use
the quotient and remainder of $x$ by $ql$.  This also covers the case
of an empty intersection.

Sum \eqref{eq:sieve-majorant} over the progression.  Expanding its
square and using \eqref{eq:crt-count} gives
\begin{align*}
 S(x;q,a,z)
 &:=\#\{1\leq n\leq x:n\equiv a\pmod q,\ \gcd(n,P)=1\}\\
 &\leq\frac{x}{q}Q(\lambda)
       +\sum_{d,e\in D}|\lambda_d\lambda_e|\\
 &\leq\frac{x}{qG}+z^2.
\end{align*}
Every prime counted by $\pi(x;q,a)$ and larger than $z$ survives this
sieve.  At most $z$ other primes occur.  This proves the first
inequality of \eqref{eq:finite-sieve}.

\subsection{A denominator bound using finite geometric sums}

Put
\[
 H_q(z)=\sum_{\substack{1\leq n\leq z\\\gcd(n,q)=1}}\frac1n.
\]
For squarefree $d$, group the integers $n\leq z$ whose radical is $d$.
Every prime exponent in such an $n$ is at most $z$, whence
\begin{equation}
 \sum_{\substack{n\leq z\\\operatorname{rad}(n)=d}}\frac1n
 \leq\prod_{p\mid d}\left(\sum_{j=1}^{z}p^{-j}\right)
 \leq\prod_{p\mid d}\frac1{p-1}
 =\frac1{\varphi(d)}.
 \label{eq:radical-bound}
\end{equation}
The second inequality uses the finite identity
$\sum_{j=1}^{z}p^{-j}=(1-p^{-z})/(p-1)$, not convergence of an
infinite Euler product.  Summing \eqref{eq:radical-bound} over all
possible radicals yields $H_q(z)\leq G_q(z)$.

Next, uniquely write $n=bc$, where all prime factors of $b$ divide $q$
and $\gcd(c,q)=1$.  Dropping the restriction $bc\leq z$ while retaining
$b,c\leq z$ gives
\begin{align*}
 H_z
 &\leq H_q(z)
       \sum_{\substack{b\leq z\\p\mid b\Rightarrow p\mid q}}\frac1b\\
 &\leq H_q(z)\prod_{p\mid q}\left(\sum_{j=0}^{z}p^{-j}\right)\\
 &\leq H_q(z)\prod_{p\mid q}\frac p{p-1}
  =\frac q{\varphi(q)}H_q(z).
\end{align*}
For $q=1$ the product is empty and equals one.  Thus
\begin{equation}
 G_q(z)\geq \frac{\varphi(q)}qH_z.
 \label{eq:sieve-denominator}
\end{equation}
This proves the second inequality of \eqref{eq:finite-sieve} and
completes the proof of Theorem~\ref{thm:finite-sieve}.

\subsection{An explicit primitive recursive threshold}

\begin{theorem}[Effective asymptotic Brun--Titchmarsh]
\label{thm:asymptotic-bt}
For every integer $k\geq1$, set $T(k)=(48k^2)^{16k}$.  If
$\gcd(a,q)=1$ and $x\geq qT(k)$, then
\begin{equation}
 \pi(x;q,a)\leq
 \left(2+\frac1k\right)\frac{x}{\varphi(q)\log(x/q)}.
 \label{eq:asymptotic-bt}
\end{equation}
The threshold is uniform in $q$ and $a$.
\end{theorem}

\begin{proof}
Let
\[
 t=x/q,\qquad v=t^{1/(16k)},\qquad
 \alpha=\frac{8k-1}{16k},\qquad z=\lfloor t^\alpha\rfloor.
\]
The hypothesis says $v\geq48k^2$.  There is no need to compute a real
power to obtain $z$: it is the largest integer $n\leq x$ satisfying
\begin{equation}
 q^{8k-1}n^{16k}\leq x^{8k-1}.
 \label{eq:root-bounded-search}
\end{equation}
This is a bounded primitive recursive search.

The elementary rectangle inequality gives $H_z\geq\log(z+1)$.
Since $z+1>t^\alpha$, the main term in
\eqref{eq:finite-sieve} is bounded by
\begin{equation}
 \frac{x}{\varphi(q)H_z}
 \leq\frac{16k}{8k-1}\frac{x}{\varphi(q)\log t}
 \leq\left(2+\frac1{3k}\right)
          \frac{x}{\varphi(q)\log t}.
 \label{eq:bt-main-term}
\end{equation}
For the last inequality, multiply positive denominators; the result
is $6k\leq8k-1$, true for $k\geq1$.

Write $V=x/(\varphi(q)\log t)$.  Since $\varphi(q)\leq q$,
$V\geq t/\log t$.  Also $z\geq1$, so
$z^2+z\leq2t^{2\alpha}=2t/v^2$.  Using
$\log v\leq v-1$ and $\log t=16k\log v$, we obtain
\begin{equation}
 \frac{z^2+z}{V}
 \leq\frac{2\log t}{v^2}
 =\frac{32k\log v}{v^2}
 \leq\frac{32k(v-1)}{v^2}
 \leq\frac{32k}{v}
 \leq\frac2{3k}.
 \label{eq:bt-error-term}
\end{equation}
Adding \eqref{eq:bt-main-term} and \eqref{eq:bt-error-term} proves
\eqref{eq:asymptotic-bt}.
\end{proof}

For a rational $\epsilon>0$, choosing
$k=\max(1,\lceil1/\epsilon\rceil)$ gives constant at most
$2+\epsilon$.  This is a fully specified asymptotic upper bound,
including the dependence of its threshold on the requested error.

\subsection{Its precise arithmetical reading}

The finite theorem \eqref{eq:finite-sieve} requires no logarithm.  To
translate \eqref{eq:asymptotic-bt}, for a rational $t\geq1$ put
$w=(t-1)/(t+1)$ and define
\[
 L_M(t)=2\sum_{j=0}^{M-1}\frac{w^{2j+1}}{2j+1}.
\]
The geometric-tail calculation gives, for $M\geq1$,
\[
 0\leq\log t-L_M(t)
 \leq\frac{2w^{2M+1}}{(2M+1)(1-w^2)}
 \leq\frac tM,
\]
where the last step uses $1-w^2\geq1/t$.  Thus
$M=2^r\lceil t\rceil$ gives error at most $2^{-r}$.
The theorem has the following first-order arithmetic form: for all
precision parameters $r$ and all $k,q,a,x$ satisfying its decidable
premises,
\begin{equation}
 k\varphi(q)\pi(x;q,a)
 L_{2^r\lceil x/q\rceil}(x/q)\leq(2k+1)x.
 \label{eq:arithmetical-abt}
\end{equation}
Every term in \eqref{eq:arithmetical-abt} is the output of a fixed
primitive recursive program.

The additional verification obligations for this formulation are the
elementary logarithm laws used above: monotonicity, the power law,
$H_z\geq\log(z+1)$, and $\log v\leq v-1$.  Their proofs must be
carried out for rational approximants with error bounds.  The positive
root used in the exposition has a primitive recursive bisection name
and may be eliminated by such approximants; the actual sieve parameter
is the bounded search \eqref{eq:root-bounded-search}.  Importing a Lean
theorem about its usual real numbers would not, by itself, discharge
these iPRA obligations.

\subsection{Why the quadratic bound needs a further argument}

At a fixed ratio $x/q$, the condition $x/q\geq T(k)$ bounds the
admissible $k$.  Consequently one cannot let $k$ tend to infinity in
Theorem~\ref{thm:asymptotic-bt} to obtain constant $2$.  The sharp
all-range theorem requires an additional argument that is absent from
the elementary remainder bound $z^2$.  Section~\ref{sec:sharp-bt}
supplies that argument using a weighted large sieve and an explicit
finite-grid estimate.

For a complete Lean development of the result proved here, the useful
order is: divisor identities; finite M\"obius inversion;
\eqref{eq:sieve-diagonal}; normalization and the weight bound;
the CRT count; the quadratic majorant; the radical/geometric
denominator estimate; and finally the certified-logarithm corollary.
Each theorem in the finite part concerns finite lists and rational
identities.  A separate compilation or direct proof-object layer must
still connect those Lean proofs to the intended iPRA proof checker.
